%BeginMC
\question
The outwelling hypothesis refers to

\begin{choices}
\correctchoice the export of excess primary production from high productivity estuaries. %EndChoice
\choice the export of larval individuals from estuary nursery grounds to the open ocean.%EndChoice
\choice the export of organic material out of the estuary by a microbial loop.%EndChoice
\choice the enhancement of offshore productivity due to larval individuals leaving low productivity estuaries.%EndChoice
\choice the loss of primary productivity to consumers in low productivity estuaries.%EndChoice
\end{choices}
%EndMC

%BeginMC
\question
Which of the following does not appear to be an important biotic controlling factor of zonation in the rocky intertidal zone?

\begin{choices}
\correctchoice Parasitism %EndChoice
\choice  Competition  %EndChoice
\choice Predation %EndChoice
\choice Grazing %EndChoice
\choice Larval recruitment %EndChoice
\end{choices}
%EndMC

%BeginMC
\question
Wind is the dominant abiotic factor determining community structure in which type of estuary?

\begin{choices}
\correctchoice Bar-built. %EndChoice
\choice Delta. %EndChoice
\choice Fjord. %EndChoice
\choice Coastal plain. %EndChoice
\choice Tectonic. %EndChoice
\end{choices}
%EndMC

%BeginMC
\question
Which two of the these five factors are the most important causes of zonation in sandy beaches?
\begin{multicols}{3}
1. Competition \\
2. Predation \\
3. Sediment particle size\\

\columnbreak

4. Space\\
5. Wave action\\
\phantom{6.} \null

\columnbreak

\phantom{7.}

\end{multicols}\vspace*{-1ex}

\begin{choices}
\correctchoice 3 and 5 %EndChoice
\choice 1 and 2 %EndChoice
\choice 1 and 3  %EndChoice
\choice 4 and 5 %EndChoice
\choice 2 and 4 %EndChoice
\end{choices}
%EndMC

%BeginMC
\question
The rocky intertidal and the rocky subtidal share one zone in common. That zone is

\begin{choices}
\choice infralittoral fringe. %EndChoice
\choice midlittoral zone. %EndChoice
\correctchoice the supralittoral. %EndChoice
\choice sublittoral fringe. %EndChoice
\choice littoral zone. %EndChoice
\end{choices}
%EndMC

%BeginMC
\question
The hypothesis proposed to explain coral reef fish diversity that assumes competition is strong and post-recruitment processes are not modified is the 

\begin{choices}
\correctchoice lottery model. %EndChoice
\choice competition model. %EndChoice
\choice predation-disturbance model. %EndChoice
\choice recruitment limitation model. %EndChoice
\choice spatial heterogeneity model. %EndChoice
\choice temporal dynamics model. %EndChoice
\end{choices}
%EndMC

%BeginMC
\question
The dynamics of patches in rocky intertidal communities is most determined by

\begin{choices}
\correctchoice patch size. %EndChoice
\choice wave action. %EndChoice
\choice ice and wrack. %EndChoice
\choice competition. %EndChoice
\choice predation. %EndChoice
\choice the species that colonizes the patch first. %EndChoice
\end{choices}
%EndMC

%BeginMC
\question
Hypersaline bare patches are most likely to occur in

\begin{choices}
\correctchoice northern salt marshes. %EndChoice
\choice any salt marsh. %EndChoice
\choice southern salt marshes. %EndChoice
\choice northern mangrove forests. %EndChoice
\choice southern mangrove forests. %EndChoice
\choice any mangrove forest. %EndChoice
\end{choices}
%EndMC
